\part{The Ambient IoT Idea}

\chapter{Why Ambient IoT?}
\label{ch:intro}

\begin{defbox}[In this chapter you will learn]
\begin{itemize}
 \item why batteries, not radios, are the real bottleneck of massive IoT;
 \item the lineage---RFID, ambient backscatter, cellular IoT---that leads to \emph{3GPP Ambient IoT};
 \item how 3GPP classifies Ambient IoT devices (A, B, C) and connects them to a network;
 \item the four properties that make these devices hard to reason about, and in particular the coupling of \emph{energy state} to \emph{communication ability};
 \item what gap Amber fills, and what questions a simulator of this kind can and cannot answer.
\end{itemize}
\end{defbox}

\section{The battery is the bottleneck}

The promise of the Internet of Things is to attach a communicating sensor to
everything worth measuring: soil moisture in a field, a pallet in a warehouse,
a strain gauge in a bridge, a temperature probe in a cold chain. Forecasts
speak of tens of billions of devices. Two properties dominate whether such a
number is reachable: the \emph{cost} of each device and the \emph{labour} of
keeping it alive.

Consider the energy budget of the humble coin cell. A typical lithium coin
cell stores about $\SI{225}{\mAh}$ at $\SI{3}{V}$, that is
\[
 E_{\mathrm{cell}} \;=\; 0.225\,\mathrm{Ah}\times 3\,\mathrm{V}\times 3600\,\mathrm{s/h}
 \;=\; \SI{2430}{J}.
\]
\begin{example}[How long does a coin cell last?]
A sensor whose \emph{average} power draw is $\SI{100}{\micro W}$ drains the cell in
\[
 t=\frac{2430\,\mathrm{J}}{100\times10^{-6}\,\mathrm{W}}=2.43\times10^{7}\,\mathrm{s}\approx 281\ \text{days}.
\]
At $\SI{10}{\micro W}$ the same cell lasts about $7.7$ years (ignoring
self-discharge, which in practice shortens this). A cellular radio that
transmits at $\SI{23}{dBm}\approx\SI{200}{mW}$, even for a fraction of a
second per day, quickly dominates the budget.
\end{example}
Now scale up. A deployment of $10^{6}$ devices with a five-year battery life
needs, on average, $10^{6}/(5\times365)\approx 548$ battery replacements \emph{per
day}, each requiring a human to find and open a device that may be inside a
wall, on a streetlight or in a sealed container. Batteries also contain
materials one would rather not scatter by the billion. The conclusion is
simple: for the very largest and cheapest deployments, the only sustainable
design is a device that needs \emph{no battery to be replaced}.

\begin{keyidea}
Massive IoT is not limited by how cheaply we can make a radio. It is limited by
how cheaply we can \emph{power} one for a decade without touching it. Ambient
IoT attacks the problem at its root: reduce the device's power to the
microwatt regime and feed it from the environment.
\end{keyidea}

\section{Three ancestors}
\label{sec:ancestors}

Ambient IoT did not appear from nowhere. It sits where three technologies
meet.

\paragraph{RFID.}
A passive UHF RFID tag contains an antenna and a tiny chip, but no battery. A
\emph{reader} radiates a continuous wave (the carrier). The tag's antenna
captures some of that power, a rectifier turns it into a DC supply for the chip,
and the chip answers by \emph{switching its antenna load} so that the
reflected carrier is modulated with the tag's data. The idea goes back to
Stockman's 1948 paper on ``communication by means of reflected power''
\citep{stockman1948}. Two features of RFID will recur throughout this book:
\emph{the reader talks first}, and \emph{the tag answers by reflecting}. RFID
readers inventory many tags using \emph{framed slotted ALOHA} (Chapter~\ref{ch:aloha}).

\paragraph{Ambient backscatter.}
In 2013 Liu et al.\ showed that devices could communicate by backscattering
\emph{ambient} signals such as TV broadcasts, with no dedicated reader at all
\citep{liu2013ambient}. The survey by Huynh et al.\
\citep{huynh2018ambient}, cited by the Amber paper, organises the field into
\emph{monostatic}, \emph{bistatic} and \emph{ambient} topologies depending on
where the carrier comes from (Chapter~\ref{ch:backscatter}).

\paragraph{Cellular IoT.}
NB-IoT and LTE-M bring standardised, wide-area, secure connectivity to
low-power devices, but they are \emph{active radios}: they generate their own
carrier, which costs tens to hundreds of milliwatts while transmitting. They
therefore need a battery.

\begin{center}
\small
\renewcommand{\arraystretch}{1.2}
\begin{tabularx}{\linewidth}{@{}>{\sffamily\bfseries}p{2.6cm}XXX@{}}
\toprule
 & Passive UHF RFID & 3GPP Ambient IoT & NB-IoT / LTE-M\\
\midrule
Energy source & RF from reader & RF, light, \dots\ (+ small storage) & Battery\\
Uplink radio & Backscatter & Backscatter (or low-power active) & Active transmitter\\
Typical active power & $\sim$\si{\micro W}--tens of \si{\micro W} & $\sim$\SI{1}{\micro W}--few hundred \si{\micro W} & $\gg$ \si{mW}\\
Who initiates & Reader & Base station / reader & Device or network\\
Network & Dedicated readers & Cellular infrastructure & Cellular infrastructure\\
\bottomrule
\end{tabularx}
\end{center}
\noindent\emph{Orders of magnitude only; the Ambient-IoT column follows the 3GPP design targets quoted below.}

Ambient IoT can be read as ``RFID's energy and backscatter physics, with
cellular IoT's infrastructure and standardisation''. That combination is what
3GPP set out to specify.

\section{3GPP's view of Ambient IoT}
\label{sec:3gpp}

3GPP studied \emph{Ambient IoT} (originally ``ambient power-enabled IoT'') in
Release~18 and Release~19 and is extending it in Release~20, with the vision
of carrying it into 6G. The Amber paper cites the overview by L\'opez et al.\
\citep{lopez2025zero}. Public 3GPP material (the study report TR~38.848, the
Release-19 solutions study TR~38.769, and the device radio specification
TS~38.191) defines the following \emph{device categories} by how much energy
they can store and whether they can generate a radio signal on their own.

\begin{center}
\small
\renewcommand{\arraystretch}{1.25}
\begin{tabularx}{\linewidth}{@{}>{\sffamily\bfseries}l>{\raggedright}X>{\raggedright}X>{\raggedright\arraybackslash}X@{}}
\toprule
Type & Energy storage & Independent signal generation & Uplink\\
\midrule
Device A & none & no & backscatter of an external carrier\\
Device B & yes (small) & no & backscatter; stored energy may be used to amplify the reflection\\
Device C & yes (small) & yes (active RF components) & active transmission\\
\bottomrule
\end{tabularx}
\end{center}
In the Release-19 normative work the focus was on the lowest-power class,
described in public summaries as a device with \emph{about
$\SI{1}{\micro W}$ peak power consumption}, some energy storage and an RF
\emph{envelope-detector} receiver, whose uplink is backscatter on a carrier
\emph{provided externally} (typically by the base station). Release~20
extends the scope towards higher-power classes. The storage in B and C is
\emph{orders of magnitude smaller} than in an NB-IoT device; this smallness is exactly why
\emph{capacitor} models matter later in this book.

\paragraph{Topologies.}
The study also considered how devices attach to the network:
\begin{enumerate}[label=\textbf{T\arabic*.}]
 \item base station $\leftrightarrow$ Ambient IoT device (direct);
 \item base station $\leftrightarrow$ \emph{intermediate node} $\leftrightarrow$ device;
 \item base station $\leftrightarrow$ \emph{assisting node} $\leftrightarrow$ device $\leftrightarrow$ base station, in which a separate node assists the downlink (for example by supplying a carrier) while the device answers the base station directly---a \emph{bistatic}-style arrangement (we describe the roles loosely; consult the report for the exact definition);
 \item user equipment $\leftrightarrow$ device.
\end{enumerate}
Amber implements topology~T1 only: base stations talk directly to nodes, and
the same base station that radiates the carrier also receives the reflection.

\begin{reality}[Where do Amber's nodes sit in this taxonomy?]
Amber's \code{Node} carries a \code{node_type} field (``passive'' or
``active''), but the field is currently \emph{cosmetic}: it only chooses a
colour when a coverage map is drawn. What really defines an Amber node is its
behaviour: it has a capacitor (so it has storage), it uplinks by
reflecting the base station's carrier (no signal generation), and it does not
amplify the reflection. That makes it closest to a storage-equipped
backscatter device (3GPP's Device~B / ``type~1''). Note, however, that the
\emph{example currents} in the shipped demos (a few mA while transmitting;
see Chapter~\ref{ch:controller}) are far above the microwatt targets of
Release~19: they model a sensor-class node, not an ultra-low-power tag. All
such numbers are configurable.
\end{reality}

\section{Four properties that make these devices hard to reason about}
\label{sec:fourprops}

It is worth stating, once and clearly, what is \emph{different} about an
Ambient IoT device from a conventional wireless node. Each property below
reappears as a design decision inside Amber.

\begin{description}[leftmargin=0.6cm,font=\sffamily\bfseries\color{ablue}]
 \item[P1 --- Energy intermittency.] The harvested power is tiny (microwatts to
   a milliwatt) and may vary. The device cannot run continuously; it
   accumulates charge, wakes up, does a little work and falls asleep again.
   Whether it is \emph{present} at a given moment is a stochastic process in
   its own right. (Chapters~\ref{ch:harvest}--\ref{ch:capacitor}.)
 \item[P2 --- Reflective uplink.] With no oscillator, mixer or power amplifier
   the device can only \emph{modulate the reflection} of someone else's signal.
   Its transmit ``power'' is a fraction of what it received, and the signal
   it creates is subject to the path loss of \emph{both} hops
   (Chapters~\ref{ch:pathloss} and~\ref{ch:backscatter}).
 \item[P3 --- Reader-driven medium access.] A device that cannot afford a
   receiver chain cannot carrier-sense, and without a carrier it cannot
   transmit at all. The base station therefore sequences everything:
   commands, reply slots, acknowledgements
   (Chapters~\ref{ch:aloha}--\ref{ch:protocols}).
 \item[P4 --- Massive scale and limited visibility.] There may be hundreds of
   devices per cell, and the base station \emph{cannot observe their energy
   state}. It sees only whether a reply arrived.
\end{description}

The Amber abstract compresses the consequence of these four properties into
one sentence that we will return to repeatedly:

\begin{keyidea}[The central coupling]
In an Ambient IoT network, \emph{a node's energy state directly determines
its ability to communicate}.
\end{keyidea}

Suppose a base station broadcasts a request and expects one hundred replies. If
sixty devices are momentarily below their activation threshold, the
protocol sees sixty ``failures'' that are \emph{not} radio failures at all.
A simulator that models radio but not energy would blame the channel; a
simulator that models energy but not protocols would not know that anything
failed. Seeing both at once is what Amber is for.

\section{The simulation gap}
\label{sec:gap}

Why not just build a testbed? The Amber paper gives the answer in its introduction:
\emph{the scale of Ambient IoT deployments, combined with unpredictable energy
availability and limited visibility into device state, makes field testing and
validation impractical.} Simulation is the natural tool, but existing ones
leave a gap, which the authors characterise as follows.
\begin{itemize}
 \item \textbf{Signal-level simulators} such as Sionna \citep{hoydis2022sionna}
   and ray-tracing platforms \citep{tarneberg2023raytracing} model the radio
   channel in great detail but, for lack of energy-harvesting, capacitor and
   state-dependent power-consumption models, do not support Ambient IoT nodes.
 \item \textbf{Network simulators} that study MAC and PHY performance
   typically assume \emph{always-available power}.
 \item \textbf{Energy-harvesting studies} usually omit protocol-level
   communication.
\end{itemize}
Amber integrates the three missing pieces in one open-source, event-driven
simulator: \emph{RF energy harvesting, capacitor storage dynamics and
backscatter communication}. The authors state it is, to their knowledge, the
first to do so. It lets a researcher ask system-level questions of the form:

\begin{quote}
\emph{If I put $N$ nodes in this cell, with this antenna, this capacitor and this
protocol, what fraction of them can I register, how often can each report, and
what is lost to collisions versus to empty capacitors?}
\end{quote}

\paragraph{What a system-level simulator is \emph{not}.}
It is not a bit-accurate PHY simulator: Amber does not simulate waveforms,
modulation, coding or fading (Chapters~\ref{ch:backscatter} and
\ref{ch:limits}). It is not hardware-validated either: the paper is a demo
abstract and reports no comparison with measurements. Treat Amber's numbers as
\emph{comparisons between designs under shared, explicit assumptions}, not as
predictions of absolute field performance. Learning to do this honestly is one
of the goals of this book.

\section{Chapter summary}

\begin{itemize}
 \item Batteries limit massive IoT; Ambient IoT removes them by combining RF/solar
   harvesting, small storage and backscatter uplinks.
 \item 3GPP defines Device~A (no storage, backscatter), B (storage,
   backscatter), C (storage, active); Release-19 targeted a $\sim\SI{1}{\micro W}$
   class.
 \item Four properties---intermittency, reflective uplink, reader-driven access,
   scale with limited visibility---couple energy to communication.
 \item Amber closes the gap between signal-level and energy-harvesting studies
   with a system-level, discrete-event model.
\end{itemize}

\section*{Exercises}
\begin{exercise}\label{ex:battery}
A $\SI{225}{\mAh}$, $\SI{3}{V}$ cell powers a sensor with an average load of
$\SI{20}{\micro W}$. Estimate the lifetime in years. By what factor must
the average load fall to reach a $15$-year lifetime?
\end{exercise}
\begin{exercise}\label{ex:swaps}
A city deploys $2\times10^{6}$ battery-powered sensors with a four-year
battery life. How many replacements per day does the municipality face
on average? Why might this figure \emph{underestimate} the labour involved?
\end{exercise}
\begin{exercise}\label{ex:classify}
A tag stores energy in a $\SI{100}{\micro F}$ capacitor, backscatters a carrier
supplied by a base station, and cannot generate any RF signal itself.
Which 3GPP device category fits best? What would change if it contained a
local oscillator and a small power amplifier?
\end{exercise}
\begin{exercise}\label{ex:reader}
Explain in three sentences why a device with no energy storage cannot send
anything when the reader's carrier is absent, and list two protocol-design
consequences of ``the reader talks first''.
\end{exercise}
